\documentclass[12pt,a4paper]{article}
\usepackage{geometry}
\geometry{margin=1in}
\usepackage{titlesec}
\usepackage{graphicx}
\usepackage{enumitem}
\usepackage{booktabs}

% Section formatting
\titleformat{\section}{\bfseries\large}{\thesection}{1em}{}
\titleformat{\subsection}{\bfseries}{\thesubsection}{1em}{}

\begin{document}

% Title Page
\begin{center}
    \Large \textbf{Capstone Project Monthly Progress Report}\\[0.5cm]
    \normalsize
    
    \textbf{Month:} March 2026 \\[0.5cm]
    \textbf{Project Title:} Uncertainty-Aware Visual Question Answering on Kvasir-VQA-x1 \\[0.5cm]
    \textbf{Project ID:} CAPSTONE\_2022\_22172\_4 \\[0.5cm]
    \textbf{Student Name:} Sunit Soni \\
    \textbf{Student Roll Number:} AP22110011494 \\

    \textbf{Total Number of Students in Team:} 03 \\
    \textbf{Details of Team Members:}\\
       Prashant Dhimal (Roll No. AP22110011492)\\
       Jayash Shrestha (Roll No. AP22110011481)\\[0.5cm]
    \textbf{Supervisor:} Dr. M Krishna Siva Prasad \\[0.5cm]
    Department of Computer Science and Engineering\\
    SRM University AP
\end{center}

\hrule
\vspace{1cm}

\section{Project Overview}
This project focuses on developing a Visual Question Answering (VQA) system for gastrointestinal medical images using the Kvasir-VQA-x1 dataset. The primary objective is to design an uncertainty-aware multimodal model that can answer clinical visual questions while minimizing hallucinations through a confidence-based abstention mechanism. The core research question investigates how effective uncertainty-aware training objectives are in reducing hallucinations in medical VQA tasks. The project targets safe and responsible deployment of multimodal AI in clinical decision support.

\section{Work Completed During Last 1 Month}
\begin{itemize}[leftmargin=*]
    \item \textbf{Comprehensive Evaluation Metrics:} Implemented a full suite of NLG evaluation metrics beyond word F1: BLEU-1/2/3/4, ROUGE-L, METEOR, BERTScore, word precision, and word recall. These metrics are standard in the medical VQA literature and enable direct comparison with published work.
    
    \item \textbf{Updated Baseline Evaluation:} Re-ran the zero-shot BLIP-2 baseline with all new metrics on 48 stratified test samples. Results: 0\% exact match, 23.4\% word F1, 19.5\% BLEU-1, 4.2\% BLEU-4, 20.8\% ROUGE-L, 21.8\% METEOR, and 25.9\% BERTScore F1. These numbers quantify the severity of hallucination in the untuned model.
    
    \item \textbf{LoRA Fine-Tuning Pipeline:} Built and executed an end-to-end fine-tuning pipeline for BLIP-2 using Low-Rank Adaptation (LoRA) with 8-bit quantization via \texttt{bitsandbytes}. LoRA was applied to the OPT language model's attention layers (\texttt{q\_proj}, \texttt{v\_proj}) with $r=16$, $\alpha=32$, and 0.1 dropout.
    
    \item \textbf{Bug Fix --- Causal LM Label Alignment:} Identified and fixed a critical training bug where the answer tokens were passed as a disconnected \texttt{labels} tensor. The corrected approach concatenates the prompt and answer as a single input sequence, with prompt tokens masked to $-100$ in labels so loss is computed only on answer tokens.
    
    \item \textbf{Fine-Tuning Results:} Trained on a stratified subset ($\sim$500 samples) for 3 epochs ($\sim$18 minutes on T4 GPU). Training loss decreased from 2.20 $\rightarrow$ 0.87 $\rightarrow$ 0.76. On 20 test samples: 5\% exact match (up from 0\%), 49.6\% word F1 (up from 23.4\%), and 50\% partial match rate (up from 5.6\%).
    
    \item \textbf{Uncertainty Estimation Module:} Implemented three complementary uncertainty estimation methods:
    \begin{itemize}
        \item \textbf{Predictive Entropy:} Computes token-level softmax entropy $H = -\sum p(t) \log p(t)$ at each generation step. High entropy indicates the model is spread across many candidate tokens.
        \item \textbf{MC Dropout:} Enables dropout at inference and runs $N=5$ forward passes. Measures lexical variance across generated answers using pairwise word F1. Inconsistent answers signal uncertainty.
        \item \textbf{Sequence Confidence:} Computes the normalized log-probability of the generated sequence. Low log-probability indicates less confident generation.
    \end{itemize}
    A combined uncertainty score weighting all three methods (0.4 entropy + 0.3 MC + 0.3 confidence) is used for abstention.
    
    \item \textbf{Abstention Mechanism:} Implemented threshold-based selective prediction --- if combined uncertainty exceeds threshold $\tau$, the model outputs ``I am not confident enough to answer this question'' instead of guessing. Threshold $\tau$ is tuned on the validation set to maximize selective accuracy while maintaining $\geq$80\% coverage (i.e., the model answers at least 80\% of questions).
    
    \item \textbf{Safety Evaluation Metrics:} Implemented four safety-specific metrics:
    \begin{itemize}
        \item \textbf{Risk-Coverage Curve:} Plots error rate vs.\ fraction of questions answered. Lower area under curve = better.
        \item \textbf{AUROC:} Measures how well uncertainty separates correct from incorrect predictions.
        \item \textbf{ECE (Expected Calibration Error):} Measures alignment between confidence and actual accuracy.
        \item \textbf{Selective Accuracy:} Accuracy at various coverage levels (50\%, 60\%, 70\%, 80\%, 90\%, 100\%).
    \end{itemize}
    
    \item \textbf{Google Colab Notebooks:} Created/updated three self-contained Colab notebooks:
    \begin{itemize}
        \item \texttt{baseline\_inference\_colab.ipynb} --- Updated with all evaluation metrics
        \item \texttt{finetune\_blip2\_colab.ipynb} --- Complete LoRA fine-tuning pipeline
        \item \texttt{uncertainty\_eval\_colab.ipynb} --- Uncertainty estimation + abstention + safety plots
    \end{itemize}
\end{itemize}

\section{Work Planned for This Month}
\begin{itemize}[leftmargin=*]
    \item Scale fine-tuning to 2,000+ stratified training samples for improved performance.
    \item Resolve remaining issues in the uncertainty evaluation notebook and run full evaluation.
    \item Analyze safety metrics (Risk-Coverage curves, AUROC, ECE) on the scaled model.
    \item Produce side-by-side comparison: baseline vs.\ fine-tuned vs.\ uncertainty-aware.
    \item Write the final project report.
\end{itemize}

\section{Progress Status}
\begin{itemize}[leftmargin=*]
    \item Overall completion: $\sim$80\% (All core modules implemented: baseline, fine-tuning, uncertainty estimation, abstention, safety metrics).
    \item Remaining: scaled training run, full uncertainty evaluation, and final documentation.
    \item Timeline is on track --- core research contributions (uncertainty estimation + abstention) are fully implemented.
\end{itemize}

\section{Challenges/Issues Faced}
\begin{itemize}[leftmargin=*]
    \item \textbf{Causal LM label alignment:} The initial fine-tuning produced garbage outputs (repeating dashes). Root cause was incorrect label construction --- the answer was passed as a disconnected tensor instead of being concatenated with the prompt in the input sequence. Debugging and fixing this consumed significant effort.
    \item \textbf{GPU resource constraints:} Limited to Google Colab T4 (16GB VRAM), requiring 8-bit quantization, small batch sizes (4), and reduced training subsets. Full dataset fine-tuning was not feasible within session time limits.
    \item \textbf{Image availability:} Only $\sim$500 of the planned 2,000 training images were available in the initial run, limiting fine-tuning effectiveness.
    \item \textbf{MC Dropout with quantized models:} Running stochastic forward passes on 8-bit quantized models required careful handling of dropout layer states to avoid interference with quantization.
\end{itemize}

\section{Learning Outcomes}
\begin{itemize}[leftmargin=*]
    \item Learned Low-Rank Adaptation (LoRA) for parameter-efficient fine-tuning of large vision-language models, including proper integration with 8-bit quantization.
    \item Understood the critical importance of label alignment in causal language model fine-tuning --- prompt tokens must be masked to $-100$ so loss applies only to answer tokens.
    \item Gained practical experience with NLG evaluation metrics (BLEU, ROUGE-L, METEOR, BERTScore) and their individual strengths for medical text evaluation.
    \item Studied and implemented three uncertainty estimation techniques (predictive entropy, MC Dropout, sequence confidence) and their complementary strengths.
    \item Learned about safety evaluation frameworks for medical AI: Risk-Coverage analysis, calibration, and selective prediction.
    \item Gained experience designing threshold-based abstention mechanisms with coverage constraints.
\end{itemize}

\section{Plan for Next Month}
\begin{itemize}[leftmargin=*]
    \item Scale fine-tuning to the full 2,000-sample stratified training subset.
    \item Complete the uncertainty evaluation pipeline and generate safety analysis plots.
    \item Produce final comparison table: baseline vs.\ fine-tuned vs.\ uncertainty-aware across all metrics.
    \item Write the final project report including methodology, experiments, results, and analysis.
    \item Expected Deliverables: Scaled model checkpoint, uncertainty analysis results, safety plots, and final report.
\end{itemize}

\section{Student Individual Contribution}

\noindent\textbf{Sunit Soni (AP22110011494):}
Designed and implemented the uncertainty estimation module with three methods: predictive entropy, MC Dropout (5-pass lexical variance), and sequence-level confidence. Built the abstention mechanism with threshold tuning targeting 80\% coverage. Implemented all four safety evaluation metrics (Risk-Coverage, AUROC, ECE, selective accuracy). Diagnosed and fixed the causal LM label alignment bug in the fine-tuning pipeline. Created the uncertainty evaluation Colab notebook with 4-panel safety visualization.

\vspace{0.4cm}

\noindent\textbf{Jayash Shrestha (AP22110011481):}
Implemented the comprehensive evaluation metrics suite: BLEU-1/2/3/4, ROUGE-L, METEOR, BERTScore, word precision, and word recall in the shared \texttt{train\_utils.py} module. Updated both the baseline and fine-tuning Colab notebooks with the new metrics. Re-ran the baseline evaluation with all metrics on 48 stratified test samples, establishing the updated benchmark numbers. Updated the project README with current results and project structure.

\vspace{0.4cm}

\noindent\textbf{Prashant Dhimal (AP22110011492):}
Led the LoRA fine-tuning implementation for BLIP-2, including 8-bit quantization setup, gradient accumulation configuration, and cosine learning rate scheduling. Trained the initial model on $\sim$500 stratified samples, achieving 49.6\% word F1 (up from 23.4\% baseline). Built the fine-tuning Colab notebook with sanity checks for label alignment verification. Configured the LoRA hyperparameters ($r=16$, $\alpha=32$, target modules) based on literature recommendations.

\section*{9. Supervisor's Observation}

\noindent\textbf{Student Completed:} \rule{3cm}{0.4pt}\% 
of the planned/assigned task during the last month

\vspace{0.5cm}

\noindent\textbf{Comments:}

\vspace{1cm}
\noindent\rule{\textwidth}{0.4pt}

\vspace{0.5cm}

\noindent\textbf{Number of times students meet with supervisor (Online or Offline):} \rule{1cm}{0.4pt} \\

\noindent\textbf{Satisfied with the Student's Work:} 
\hspace{0.5cm} Satisfied 
\hspace{1cm} Not Satisfied

\vspace{0.5cm}

\vspace{1.5cm}

\noindent\begin{tabular}{p{7cm} p{6cm}}
\textbf{Supervisor Signature:} & \rule{5cm}{0.4pt} \\
\\
\textbf{Date:} 
\end{tabular}

\vspace{0.8cm}

\end{document}
